\documentclass[border=1pt]{standalone}

\usepackage{tikz}
\usetikzlibrary{arrows.meta, calc}

\usepackage{luatexja} 
\usepackage{amsmath}

\begin{document}
    \begin{tikzpicture}[>=Stealth, thick]
        % パラメータの定義
        \pgfmathsetmacro{\R}{4.1}
        \pgfmathsetmacro{\Rz}{5.9}

        \pgfmathsetmacro{\rd}{2.1}
        
        \pgfmathsetmacro{\xlens}{2.6} 
        \pgfmathsetmacro{\ylensa}{{\R - sqrt(\R^2 - (\xlens)^2)}}
        \pgfmathsetmacro{\ylensb}{{\Rz - sqrt(\Rz^2 - (\xlens)^2)}}

        \pgfmathsetmacro{\yra}{\R - sqrt(\R^2 - \rd^2)}
        \pgfmathsetmacro{\yrb}{\Rz - sqrt(\Rz^2 - \rd^2)}

        \pgfmathsetmacro{\arrowgap}{\xlens/1.7}
        \pgfmathsetmacro{\larrowa}{1.2}

        \pgfmathsetmacro{\airmargin}{1}
        \pgfmathsetmacro{\larrowb}{1.4*(\yra-\yrb)}

        % 座標の定義
        \coordinate (C) at (0, 0);             
        \coordinate (Or) at (0, \R);           
        \coordinate (Orz) at (0, \Rz);         

        \coordinate (Pup) at (\rd, \yra);
        \coordinate (Plow) at (\rd, \yrb);

        \coordinate (Ea) at (\rd, \R);
        \coordinate (Eb) at (\rd+0.23, \R-0.3);
        \coordinate (Ec) at (\rd-0.1, \R-0.3);

        % 平凸レンズ
        \draw[domain=-\xlens:\xlens, samples=200] plot (\x, {\R - sqrt(\R^2 - (\x)^2)});
        \draw (-\xlens, \ylensa) -- (\xlens, \ylensa);
        \node[left] at (-\xlens, \ylensa) {\small 平凸レンズ};
        \fill (Pup) circle (1.5pt);

        % 平凹レンズ
        \draw[domain=-\xlens:\xlens, samples=200] plot (\x, {\Rz - sqrt(\Rz^2 - (\x)^2)});
        \draw (-\xlens, \ylensb) --++ (0,-\ylensb-0.5*\ylensa)
            --++ (2*\xlens,0) --++ (0,\ylensb+0.5*\ylensa);
        \node[left] at (-\xlens, 0) {\small 平凹レンズ};
        \fill (Plow) circle (1.5pt);

        % R,R0,r'の寸法線
        \draw[semithick] (C) to[bend left=18] node[fill=white] {$R$} (Or); 
        \draw[semithick] (C) to[bend right=12] node[fill=white] {$R_0$} (Orz);
        \draw[<->,semithick] (0, -0.24*\ylensa) --++ (\rd,0)
            node[midway, fill=white, yshift=1pt, inner sep=2pt] {$r'$};
        \draw[semithick] (Pup) -- (\rd, -0.32*\ylensa);

        % 垂直軸
        \draw[dash pattern=on 5pt off 2pt on 2pt off 2pt, semithick]
            (0, -0.32*\ylensa) -- (0, \Rz);

        % 半径の線と中心点
        \fill (Or) circle (1.5pt);
        \draw[thin] (Or) --++ (220:0.7*\larrowa)
            node[left, inner sep=1pt, align=left, font=\setlength{\baselineskip}{12pt},font=\small]
            {平凸レンズの \\ 球面の中心};
        \fill (Orz) circle (1.5pt);
        \draw[thin] (Orz) --++ (220:0.7*\larrowa)
            node[left, inner sep=1pt, align=left, font=\setlength{\baselineskip}{12pt},font=\small]
            {平凹レンズの \\ 球面の中心};
        \draw (Or) -- (Pup);
        \draw (Orz) -- (Plow);

        % 接点 C'
        \node[below left, inner sep=2pt] at (C) {$\mathrm{C}'$};
        \fill (C) circle (1.5pt);

        % 空気層の厚さ
        \draw[semithick]
            (Pup) --++ (\airmargin, 0) (Plow) --++ (\airmargin, 0);
        \draw[<-, semithick] (Pup) ++ (0.8*\airmargin, 0) --++ (0, \larrowb);
        \draw[<-, semithick] (Plow) ++ (0.8*\airmargin, 0) --++ (0, -\larrowb);
        \node[right, align=left, font=\setlength{\baselineskip}{12pt},font=\small]
            at (\rd+\airmargin, \yra/2+\yrb/2){空気層 \\ の厚さ};

        % 単色光の矢印
        \node[above] at (0, 1.05*\Rz+0.9*\larrowa) {\small 単色光};
        \foreach \x in {-1.5*\arrowgap, -0.5*\arrowgap, 0.5*\arrowgap, 1.5*\arrowgap}
            {
                \draw[thick, <-] (\x, 1.05*\Rz) --++ (0, \larrowa);
            }
        
        % 観測者
        \draw (Ea) to[bend left=25] (Eb);
        \draw (Ea) to[bend left=14] (Ec);
        \filldraw ($(Eb)!0.5!(Ec)$) circle [x radius=0.14, y radius=0.05];
        \draw (Ec) to[bend left=20] ++(-0.04, -0.09);
        \draw (Ec) ++ (0.02, 0.01) to[bend left=20] ++(-0.08, -0.02);
        \node[right, yshift=3pt] at (Eb) {\small 観測者};

        % 右側余白の調整
        \node [right] at (\xlens, 0) {\small 　　　　　};

    \end{tikzpicture}
\end{document}